\chapter{Gestion des écarts et adaptation}
\label{chap:adaptation-aer}

\section{Qualification d'un écart}

Une adaptation commence par la comparaison d'un état attendu avec un état observé. Pour l'approvisionnement, le modèle conserve le délai nominal utilisé par la planification et calcule séparément la date effective de réception. Lorsque les deux divergent dans les conditions du scénario, l'écart devient un événement métier identifiable. Il n'est pas déduit a posteriori d'un indicateur global.

Le fournisseur ou l'agent Source signale d'abord le fait local. Le pilotage opérationnel le qualifie comme exception, le coordinateur l'inscrit dans le contexte du processus, puis le tactique révise les informations nécessaires aux domaines aval. Chaque message apporte une transformation précise; répéter la même alerte à plusieurs niveaux ne suffirait pas à démontrer une adaptation.

Le modèle permet également d'activer une panne, des perturbations de demande et des scénarios de qualité. Ces capacités restent distinctes du retard fournisseur. Elles ne reçoivent pas de résultat quantitatif lorsqu'aucune exécution primaire dédiée n'est disponible.

\section{Retard fournisseur contrôlé}

Le scénario perturbé cible le GPL fourni par \Agent{ACT\_1}. Le délai nominal vaut 6 h et le délai effectif 12 h, soit un retard additionnel de 6 h. Une seule réception est ciblée. La commande porte 20 unités et l'ordre interne associé est \Agent{REAPPRO\_1}; le besoin GPL correspondant vaut 250 unités.

Les exports Excel et ABox décrivent la même exécution et la même clôture à T=2 343 s. Ils confirment la matière, le fournisseur, les délais nominal et effectif, ainsi que la chaîne de messages. Cette concordance établit l'activation du retard et sa propagation. Elle ne mesure pas son effet causal net sur le délai final, car aucune exécution de référence sans retard n'est disponible.

La commande est livrée en totalité et en retard. Ces statuts ne permettent pas d'attribuer toute la tardiveté au fournisseur. Source, Make et Deliver possèdent leurs propres dynamiques, et la conversion temporelle doit être prise en compte avant toute décomposition quantitative.

\section{Propagation vers Make et Deliver}

La propagation AER comporte six transformations structurées. Elles relient l'écart local de Source à une révision du plan Deliver, sans modifier directement un flux physique.

\begin{table}[H]
  \centering
  \caption{Propagation informationnelle du retard fournisseur.}
  \label{tab:propagation-retard}
  \begin{tabularx}{\textwidth}{@{}L{0.25\textwidth}L{0.24\textwidth}Y@{}}
    \toprule
    Message & Passage principal & Transformation métier \\
    \midrule
    \texttt{SupplierDelayAlert} & Fournisseur vers pilotage Source & Rend visible l'écart entre réception nominale et réception effective. \\
    \shortstack[l]{\texttt{Operational}\\\texttt{Exception}} & Pilotage vers coordination Source & Qualifie l'alerte comme incident opérationnel du processus Source. \\
    \shortstack[l]{\texttt{ProcessDeviation}\\\texttt{Report}} & Coordination vers tactique Source & Consolide l'impact de l'incident au niveau du plan Source. \\
    \shortstack[l]{\texttt{RevisedMaterial}\\\texttt{Availability}} & \Agent{AT-sP2} vers \Agent{AT-sP3} & Transmet à Make une disponibilité prévisionnelle révisée. \\
    \shortstack[l]{\texttt{RevisedProduction}\\\texttt{CompletionDate}} & \Agent{AT-sP3} vers \Agent{AT-sP4} & Traduit l'information matière en révision potentielle de la fin de production. \\
    \texttt{RevisedDeliveryPlan} & \Agent{AT-sP4} vers \Agent{CA-sD1} & Demande la réévaluation du plan Deliver dans le nouveau contexte. \\
    \bottomrule
  \end{tabularx}
\end{table}

La séquence montre une adaptation structurée, sans démontrer une optimisation globale. Source consolide un fait fournisseur; Make reçoit une information révisée sur sa contrainte amont; Deliver reçoit une révision de plan. La trace ne montre pas qu'une solution optimale est recherchée parmi toutes les stratégies possibles.

\FigureOrPlaceholder{figures/propagation_retard_fournisseur.png}{Propagation informationnelle d'un retard fournisseur depuis Source vers Make puis Deliver.}{fig:propagation-retard-fournisseur}

\section{Blackboard, goulot et arbitrage local}

L'AER structure les échanges par le type du message, l'émetteur, le destinataire, le contexte de processus et l'objet concerné. Le Blackboard remplit une autre fonction: il conserve les observations et décisions qui doivent rester accessibles à plusieurs agents.

L'exécution étudiée contient une détection de goulot sur \Agent{sM1.3.1}, une décision \texttt{REBALANCE}, un score AHP local de 0,730 et les accusés d'application. Ces éléments montrent qu'un état opérationnel peut être partagé, qualifié puis transformé en instruction locale. Ils ne suffisent pas à mesurer l'effet de cette décision sur les temps ou sur le PI.

Une boucle adaptative complète exige une décision contextualisée, son application et un accusé. Le mécanisme de panne peut alimenter une boucle comparable, mais aucune panne n'est observée dans le scénario retenu.

\section{Distinction entre arbitrage AHP et Performance Index}

L'AHP et le PI répondent à deux questions différentes. L'AHP classe des options dans une situation locale ou tactique. Son score dépend des critères et de la matrice de comparaison utilisés pour cet arbitrage. Le PI évalue l'exécution à partir de mesures transformées en métriques, puis en scores RL, RS, AG, CO et AM.

Un score AHP ne constitue donc pas une composante directe du PI. L'exécution simultanée d'une décision \texttt{REBALANCE} et du pipeline de performance ne démontre pas que l'une améliore l'autre. Une telle relation demanderait des exécutions comparables avec et sans la décision étudiée.
